\section*{Capítulo \ref{ch:g7:chemical-reactions} --- As reações químicas: reagentes e produtos}

\begin{solution}{exo:g7:chemical-reactions:1}
Físicas: o gelo derretendo, o açúcar se dissolvendo. Químicas: um
fósforo queimando, o pão torrando, um prego enferrujando.
\end{solution}

\begin{solution}{exo:g7:chemical-reactions:2}
Os \omterm{def:g7:chemical-reactions:reactant}{reagentes} são as espécies consumidas; os \omterm{def:g7:chemical-reactions:reactant}{produtos} são as espécies
novas formadas.
\end{solution}

\begin{solution}{exo:g7:chemical-reactions:3}
carbono $+$ oxigênio $\longrightarrow$ dióxido de carbono.
\end{solution}

\begin{solution}{exo:g7:chemical-reactions:4}
\omterm{def:g7:chemical-reactions:reactant}{Reagentes}: o ferro e o oxigênio. \omterm{def:g7:chemical-reactions:reactant}{Produto}: o óxido de ferro.
\end{solution}

\begin{solution}{exo:g7:chemical-reactions:5}
O embaçado mostra a água; a água de cal leitosa mostra o dióxido de carbono.
\end{solution}

\begin{solution}{exo:g7:chemical-reactions:6}
zinco $+$ ácido clorídrico $\longrightarrow$ cloreto de zinco $+$ hidrogênio.
\end{solution}

\begin{solution}{exo:g7:chemical-reactions:7}
Antes: 4 (na \omterm{def:g7:atoms-and-molecules:molecule}{molécula} de metano). Depois: $2 + 2 = 4$ (dois em cada
\omterm{def:g7:atoms-and-molecules:molecule}{molécula} de água).
\end{solution}

\begin{solution}{exo:g7:chemical-reactions:8}
Não aparece nenhuma espécie nova: o sal continua lá, espalhado pela
água, e é recuperado fazendo a água evaporar. É uma \omterm{def:g7:chemical-reactions:physical-transformation}{transformação
física}.
\end{solution}

\begin{solution}{exo:g7:chemical-reactions:9}
A maior parte da madeira reagiu com o oxigênio e deu \omterm{def:g7:chemical-reactions:reactant}{produtos} que são
gases, o dióxido de carbono e o vapor de água, que vão embora para o \omterm{def:g4:air-a-mixture-of-gases:air}{ar}.
Só a cinza fica.
\end{solution}

\begin{solution}{exo:g7:chemical-reactions:10}
hidrogênio $+$ oxigênio $\longrightarrow$ água. O \omterm{def:g7:chemical-reactions:reactant}{produto} deixa azul o
sulfato de cobre anidro.
\end{solution}

\begin{solution}{exo:g7:chemical-reactions:11}
Antes: $2 \times 2 = 4$ \omterm{def:g7:atoms-and-molecules:atom}{átomos} de hidrogênio e 2 \omterm{def:g7:atoms-and-molecules:atom}{átomos} de oxigênio.
Depois: duas \omterm{def:g7:atoms-and-molecules:molecule}{moléculas} de água contêm $2 \times 2 = 4$ \omterm{def:g7:atoms-and-molecules:atom}{átomos} de
hidrogênio e $2 \times 1 =
2$ \omterm{def:g7:atoms-and-molecules:atom}{átomos} de oxigênio. Os mesmos \omterm{def:g7:atoms-and-molecules:atom}{átomos}, nos mesmos números.
\end{solution}

\begin{solution}{exo:g7:chemical-reactions:12}
Um \omterm{def:g7:atoms-and-molecules:atom}{átomo} de oxigênio não pode surgir do nada: o desenho tem 2 \omterm{def:g7:atoms-and-molecules:atom}{átomos} de
oxigênio antes e 3 depois. O correto: um \omterm{def:g7:atoms-and-molecules:atom}{átomo} de carbono $+$ uma
\omterm{def:g7:atoms-and-molecules:molecule}{molécula} de oxigênio
$\longrightarrow$ uma \omterm{def:g7:atoms-and-molecules:molecule}{molécula} de dióxido de carbono, sem nada sobrando.
\end{solution}

\begin{solution}{pb:g7:chemical-reactions:1}
\textbf{1.} O metano e o oxigênio.

\textbf{2.} A água.

\textbf{3.} O dióxido de carbono.

\textbf{4.} Não: ele não é consumido, e não é escrito na \omterm{def:g7:chemical-reactions:word-equation}{equação por
extenso}.

\textbf{5.} metano $+$ oxigênio $\longrightarrow$ dióxido de carbono $+$
água.

\textbf{6.} Química: duas espécies desaparecem (o metano e o oxigênio),
e duas espécies novas aparecem, detectadas por testes; além disso, há
liberação de calor e de luz.

\textbf{7.} \ce{CH4}, \ce{O2}, \ce{CO2}, \ce{H2O}.

\textbf{8.} 1 de carbono, 4 de hidrogênio, $2 \times 2 = 4$ de oxigênio.

\textbf{9.} 1 de carbono; $2 \times 2 = 4$ de hidrogênio;
$2 + 2 \times 1 = 4$ de oxigênio. Os mesmos \omterm{def:g7:atoms-and-molecules:atom}{átomos}, nos mesmos números: eles só foram
rearranjados.

\textbf{10.} $10 \times 2 = 20$ \omterm{def:g7:atoms-and-molecules:molecule}{moléculas} de oxigênio.

\textbf{11.} 10 \omterm{def:g7:atoms-and-molecules:molecule}{moléculas} de dióxido de carbono.

\textbf{12.} $10 \times 2 = 20$ \omterm{def:g7:atoms-and-molecules:molecule}{moléculas} de água.
\end{solution}
